\chapter{数据收集与测量}

本章主要对实证研究设计以及量表的开发进行具体操作的说明。本研究首先明确数据采集的技术路线、目标群体选择的标准和统计学特征，突出样本分布具有典型性和科学性；接着对各个研究变量的测验工具来源、题型设计原则、信效度检验方法做详细的说明，为后面实证分析提供理论基础和技术支持。可靠性与有效性测试结果说明，所建立的量表具有较好的严谨性，并有较强的实践应用性。

\section{数据收集及样本介绍}

本研究采用问卷调查法收集实证数据，把问卷设计成五个部分，即控制变量、创新创业教育模式、自主动机、外部激励、创业意愿。选取天津市某高校创业学院学生为研究对象的原因有三点，第一，该学院属于区域范围内重要的创新创业教育平台，拥有较完备的创新创业教学课程体系；第二，该群体具有较高的学历水平和相似的年龄结构，适合用来检验研究结论的有效性；第三，被调查者对问卷中的问题比较熟悉，从而提高数据的质量以及科学性。

本研究的问卷设计过程中有文献综述、量表选择和校本化改造这三个主要部分。研究人员经过查阅相关的学术文献，从中提取出与研究主题相对应的量表，然后对提取出的量表做针对性的修改改进。在正式实行之前，选32名天津某高校学生做预试样本，目的是检验题目是否准确，是否连贯，有无效度隐患。根据调查得到的建议，对问题的表达和结构做出必要的调整和完善，形成了一套适合大学生创新创业教育研究使用的标准调查问卷。为了保证数据的一致性、可靠性，所有的题项都用李克特五级量表\nolinebreak[3]（1至5分分别对应强烈反对、反对、稍微反对、支持、支持）来填写。

本研究使用腾讯问卷平台，利用IP地址唯一的标识机制进行在线调查，用创业学院课程微信群发送带有编码链接的问卷。初步结果表明，共收回有效问卷168份。经过数据清洗过程，剔除掉逻辑校验出错或者作答时间明显偏离常态的无效样本之后，得到有效的问卷有154份，有效回收率是91.67\%。

本次调查问卷的回收情况较好，样本具有较好的代表性。性别比例为49.35\%，50.65\%平衡，研究对象群体男性、女性人数大致相当。从年级来看，调查样本主要是本科层次，大一、大二年级学生各占26.62\%、25.32\%，大三学生占比较低\nolinebreak[3]（13.64\%），大四及以上学历\nolinebreak[3]（含研究生）群体所占比例为39.46\%，其中有24.68\%的大四学生和9.74\%的高阶学历者。按专业领域分为工学、经济学和管理学三个学科的比重分别为41.56\%、30.52\%、8.44\%。就创业实践经历来说，有27.27\%的受访者有过实际创业经历，剩下72.73\%的被调查对象没有参加过类似的事情；八成以上(85.71\%)的研究对象具有过志愿服务的经历。就课程偏好而言，偏好理论教学的学生占了44.16\%，偏好实践导向课程的占34.42\%，同时推崇两者相融并行混合式学习的学生有21.43\%。

\section{变量测量}

\subsection{自变量测量}

本研究把创新创业教育模式作为核心变量，将它分为理论型、实践型和融合型三种类型，并使用标准量表进行测量。在理论型课程上，评价标准依据目前有关理论导向型创业教育的研究成果来设置四个问题，即课程内容展现完整的创业理论框架以及经典案例解析明显提高个人的理解能力。对实践型课程来说，采用体验式学习理论来创建评价体系，设立四个题项，即\nolinebreak[3]“教学设计能保障实践活动的足够数量和质量”、\nolinebreak[3]“能用专业知识解决实际的创业问题”。融合型课程的指标由研究者自行设计，有四个体现综合素养改善的目标表述，比如\nolinebreak[3]“课程成功做到理论和实践的有机结合”。所有的问卷都是用李克特五级量表的形式进行填写，得分越高就表示受访者对本课程的认同度越高。

\subsection{中介变量测量}

本研究以Deci与Ryan\nolinebreak[3]（1985）提出的自我决定理论为依托，利用Vallerand等人的学术动机量表进行建构\cite{Vallerand1992}。就内部动机以及认同调节这两个主要方面而言，创建了一系列的评估指标。具体的题项有，对创业实践活动有浓厚兴趣、参与创业课程使我产生了愉悦感和成就感、期望通过学习创业课程来提高自我效能感并且实现个人价值、我认为掌握了创业技能会对今后的职业发展产生重大影响等陈述性试题。另外加入与人格特征形成相关的表述为解释变量。使用五级Likert评分法来收集数据，总得分越高，说明被试个体自主性动机水平越高。

\subsection{调节变量测量}

本研究以Deci和Ryan(1985)的自我决定理论为依据，根据Vallerand等人的学术动机分类体系，从外在动机的角度来设计测量工具\cite{Vallerand1992}。根据内摄型和外源型动机的主要内容，从五个方面来设计一个包含十个题目内容的问卷框架。本问卷包含规避负罪感参加创业课、展现个人能力选课、获得学分完成学业目标、受到奖励或者认证行为驱动四个方面的内容。评分规则说明总分越高，则个体的外部动机特征就越大。

\subsection{因变量测量}

本研究以创业意愿为关键自变量，运用Liñán和Chen\nolinebreak[3]（2009）基于计划行为理论开发的问卷工具——创业意向量表\nolinebreak[3]（Entrepreneurial Intention Questionnaire，EIQ）\cite{Linan2009}。量表中有7个代表性题项，即\nolinebreak[3]“未来参与自主创业具有很强的吸引力”，或者\nolinebreak[3]“成为创业者是个人职业发展最主要的需求”。从研究结果可以看出，在多元文化背景之下，本量表的信效度较好，得分和个体创业倾向正相关，分数越高说明其创业意愿就越强。

\subsection{控制变量}

本研究以创业领域相关文献为依托，用性别、年级、专业背景、创业经历、志愿服务这五个控制变量控制了个体特征给研究结论造成的干扰。实证结果表明性别差异具有明显的异质性特点，在风险偏好以及创业认知上都有体现\nolinebreak[3]（Zhao等，2015；Noguera等，2024）\cite{Zhao2015,Noguera2024}；年级变量体现的是学习者随着时间的推移所具有的知识储备和社会经验的阶段变化情况；学业专业方向对创业机会识别能力、资源协同效率起到着引导的作用。创业实践经验体现的是受试者经过实践阶段后创新思维以及职业素质的形成程度，而志愿服务的经历是从另外一个方面体现出一个人的社会责任感意识以及合作精神。采用上述的控制变量之后，本研究所研究出的大学生创新创业课对大学生创业意向的影响机制更加精准，可以提高本研究结果的内部一致性以及科学性、可信度。

\section{问卷信度分析}

本研究采用SPSSAU统计软件对问卷量表进行信效度检验，用Cronbach's α系数来评价量表的可靠性和稳定性。按公认的标准，α值大于0.7就表明量表的内在一致性较高。根据表\ref{tab:reliability}可知理论型课程、实践型课程、融合型课程量表的α系数分别是0.878、0.921和0.889；自主动机量表、外部动机量表的α系数分别为0.907和0.815；创业意愿量表的α系数最高，为0.964。所有量表的α系数都大于0.7的标准，说明本研究所设计的测量工具内部一致性很好。

\tablemacro{htbp}
            {各量表 Cronbach's α 系数、平均方差提取量、组合信度}
            {ccccc}
            {\songtibf{潜变量} & \songtibf{题项} & \songtibf{Cronbach's α} & \songtibf{AVE} & \songtibf{CR} \\}
            {理论型课程 & 4 & 0.878 & 0.514 & 0.920 \\
            实践型课程 & 4 & 0.921 & 0.625 & 0.908 \\
            融合型课程 & 4 & 0.889 & 0.552 & 0.876 \\
            自主动机 & 5 & 0.907 & 0.526 & 0.869 \\
            外部动机 & 5 & 0.815 & 0.653 & 0.849 \\
            创业意愿 & 7 & 0.964 & --- & --- \\}
            {tab:reliability}

\section{问卷效度分析}

本本研究用探索性因子分析的方法来检验问卷结构效度。数据分析结果表明，Kaiser-Meyer-Olkin (KMO) 取样适应性系数为0.920，大于0.7的标准，说明该数据适合做因子分析。采用主成分分析法和最大方差旋转得到5个因素，累计解释变异比例达到84.675\%，远远超过60\%的理论要求。各个题项在对应的因子上的载荷均值均大于0.5，共同度均大于0.4，以上指标充分显示出了本调查问卷的结构效度很好。

为进一步验证问卷的因素结构，本研究还进行了验证性因子分析，比较了五因子模型与多个备选模型的拟合优度，结果如表\ref{tab:validity}所示。相比之下，四因子模型的拟合效果有所下降\nolinebreak[3]（CFI = 0.96，RMSEA = 0.08），而三因子、双因子和单因子模型的拟合效果显著变差\nolinebreak[3]（CFI均低于0.50，RMSEA均高于0.30）。这一结果进一步支持了本问卷的五因子结构，表明问卷具有较好的结构效度。

\tablemacro{htbp}
            {问卷效度检验}
            {lcccccccc}
            {\songtibf{模型} & \songtibf{$\chi^2$} & \songtibf{df} & \songtibf{$\Delta\chi^2$} & \songtibf{$\Delta$df} & \songtibf{CFI} & \songtibf{RMSEA} & \songtibf{SRMR} & \songtibf{$\chi^2$/df} \\}
            {五因子模型 & 402.89 & 265.00 & --- & --- & 0.98 & 0.06 & 0.05 & 2.79 \\
            四因子模型 & 506.69 & 246.00 & 105.80 & -19.00 & 0.96 & 0.08 & 0.04 & 3.65 \\
            三因子模型 & 4041.27 & 249.00 & 3638.38 & -16.00 & 0.50 & 0.31 & 0.21 & 4.76 \\
            双因子模型 & 4161.43 & 251.00 & 3758.54 & -14.00 & 0.49 & 0.32 & 0.21 & 5.13 \\
            单因子模型 & 5867.40 & 252.00 & 5464.50 & -13.00 & 0.27 & 0.38 & 0.38 & 5.54 \\}
            {tab:validity}

\noindent 注：所有数值都保留了两位小数。

\clearpage
